\documentclass[10pt,a4paper]{amsart}
\usepackage[margin=19mm]{geometry}
\usepackage{amsmath,amssymb,mathtools}
\usepackage[scr=rsfs,cal=euler]{mathalfa}
\usepackage{xfrac}
\usepackage[unicode,hidelinks,bookmarks=false]{hyperref}
\newcommand{\ie}{i.e.\ }
\newcommand{\cf}{cf.\ }
\newcommand{\resp}{respectively\ }
\newcommand{\Subsection}{\subsection}
\providecommand{\nobreakdash}{\nobreak\hbox{-}}
\setlength{\emergencystretch}{2em}
\multlinegap0pt
\usepackage{mathtools,mathrsfs,enumitem,booktabs}
\newtheorem{theorem}{Theorem}
\newtheorem{lemma}{Lemma}
\newtheorem{corollary}{Corollary}
\newtheorem{proposition}{Proposition}
\newcommand{\Cl}{\operatorname{Cl}}
\newcommand{\atimes}{\dot{\otimes}}
\newcommand{\myd}{d}
\begin{document}
\title[Failure of Strong Convergence of Matrices with Fermionic Ent]{Failure of Strong Convergence of Matrices with Fermionic Entries}
\author{Dimitri Shlyakhtenko ạte theorem Theorem lemma Lemma corollary Corollary proposition Proposition Cl ọt d e̱gin document; Dimitri Shlyakhtenko}
\date{}
\maketitle
\noindent\emph{Source and licence.} Failure of Strong Convergence of Matrices with Fermionic Entries. Version arXiv:2606.28648v1 (CC BY 4.0). This material is distributed under the Creative Commons Attribution 4.0 International licence, http://creativecommons.org/licenses/by/4.0/, as recorded in the arXiv OAI licence metadata for this exact version and shown on the abstract page https://arxiv.org/abs/2606.28648v1. Full rights notice: licenses/arxiv-2606.28648-CC-BY-4.0.txt.\par\smallskip
\noindent\textit{Editorial scope.} Counted unit: the original \texttt{theorem} environment \texttt{thrm:noStrong} at source lines 610--617 together with its complete proof at lines 618--622; the delivered document keeps source lines 433--623 so that every referenced label and every citation resolves inside it. Native editable source is the paper's own concatenated TeX; the AMS wrapper is editorial formatting only. No publisher class, upstream script or unrelated artwork is redistributed.
\section{Failure of strong convergence}
In this section we show that already for linear combinations of words of length $2$ in the matrices $A_N^{(r)}, B_N^{(r)}$ convergence of operator norms may fail.

We start with a few preliminary computations. Fix $M,N\geq 1$ and consider the operators $Q_N^{(k)}$, $k=1,\dots,M$, acting on the tensor product $\mathbb{C}^N\otimes \mathscr{F}$ of $\mathbb{C}^N$ and the antisymmetric Fock space $\mathscr{F}$. 

\subsection{The vector $\Omega_{J,M}$}
Let \[
d^{(k)}_{pq}:=\frac12\bigl(b(h^{(1,k)}_{pq})+i b(h^{(2,k)}_{pq})\bigr)
\] so that 
\[ Q_N^{(k)} = \frac{2}{\sqrt{N}}\sum_{p,q} e_{pq} \otimes  d^{(k)}_{pq}.\]

Let 
\[
    \Omega_{J,M}
    =
    2^{-MN^2/2}
    \prod_{k=1}^M\prod_{p,q=1}^N
    \left(
      1-i\,a(h^{(1,k)}_{pq})^*a(h^{(2,k)}_{pq})^*
    \right)\Omega,
\]
The order of the product defining $\Omega_{J,M}$ does not matter, since the operators $a(h^{(1,k)}_{pq})^*$ anticommute, and thus the operators $a(h^{(1,k)}_{pq})^*a(h^{(2,k)}_{pq})^*$ commute, for different triples $(p,q,k)$.

The next Lemma shows that the vector $\Omega_{J,M}$ behaves like a vacuum vector for the operators $d_{pq}^{(k)}$, which themselves turn out to form a CAR family.

\begin{lemma} \label{lemma:propertiesOfOmega}
    With the above notation, the following properties hold: 
    \begin{enumerate}[label={\rm(\roman*)}]
    \item $\Omega_{J,M}$ is a unit vector 
    \item For every \(k,p,q\), \(
    d^{(k)}_{pq}\Omega_{J,M}=0. \) 
    \item For every \(k,p,q,r\), \(
    d^{(k)}_{rq}(d^{(k)}_{pq})^*\Omega_{J,M}
    =
    \delta_{rp}\Omega_{J,M}.
    \)
    \end{enumerate}
\end{lemma}
\begin{proof}
Let \[\Omega_{J, p,q,k}= \frac{1}{\sqrt{2}}\left(
      1-i\,a(h^{(1,k)}_{pq})^*a(h^{(2,k)}_{pq})^*
    \right)\Omega = \frac{1}{\sqrt{2}}\left(\Omega - i h^{(1,k)}_{pq}\wedge h^{(2,k)}_{pq}\right).\] It is clear that 
\(\Omega_{J, p,q,k}\) is a unit vector, and is an even-parity vector in the 
Fock space. 
Under the canonical graded tensor product decomposition of the Fock space
over the orthogonal direct sum of the two-dimensional spaces
\[
\operatorname{span}\{h_{pq}^{(1,k)},h_{pq}^{(2,k)}\},
\]
the vector \(\Omega_{J,M}\) is the tensor product of the unit even-parity vectors
\(\Omega_{J,p,q,k}\).  Since the tensor products of orthogonal unit vectors are unit vectors, we obtain (i).

For part (ii), note that $d^{(k)}_{pq}$ commutes with $1-i\,a(h^{(1,k')}_{p'q'})^*a(h^{(2,k')}_{p'q'})^*$ unless $p=p'$, $q=q'$, $k=k'$.  It follows that 
\[d^{(k)}_{pq}\Omega_{J,M} = \left(d^{(k)}_{pq} \Omega_{J,p,q,k}\right) 
\wedge \left(\bigwedge_{(p',q',k')\neq (p,q,k)} 
\Omega_{J,p',q',k'}\right), \] where the order of the antisymmetric tensor products is again irrelevant since all but one of the terms have even parity.

But
\begin{eqnarray*}
    2\sqrt{2} d^{(k)}_{pq} \Omega_{J,p,q,k} &=& 
    \bigl(b(h^{(1,k)}_{pq})+i b(h^{(2,k)}_{pq})\bigr) \bigl(\Omega - i h^{(1,k)}_{pq}\wedge h^{(2,k)}_{pq}\bigr) \\ &= & 
    h^{(1,k)}_{pq} + i h^{(2,k)}_{pq} \\
    && - h_{pq}^{(1,k)}\wedge h^{(1,k)}_{pq}\wedge h^{(2,k)}_{pq} - i h_{pq}^{(2,k)} \\
    &&  - i h_{pq}^{(2,k)} \wedge h^{(1,k)}_{pq}\wedge h^{(2,k)}_{pq}
     - h_{pq}^{(1,k)}  \\ &=& 0.
\end{eqnarray*} Therefore, (ii) holds. 

For part (iii), note that 
\[
d_{pq}^{(k)}
=
\frac12\Big(
a(h^{(1,k)}_{pq})
+a(h^{(1,k)}_{pq})^*
+i\,a(h^{(2,k)}_{pq})
+i\,a(h^{(2,k)}_{pq})^*
\Big).
\] 

We claim that 
\begin{equation} \label{eqn:CARd}
\{(d^{(k)}_{rq}), (d^{(k)}_{pq})^*\} = \delta_{rp} 1.
\end{equation}

If $r\neq p$, the operators $d_{rq}^{(k)}$ and $(d_{pq}^{(k)})^*$ anti-commute, since the corresponding expressions in the creation/annihilation operators $a^*$ and $a$ anti-commute, and thus \eqref{eqn:CARd} holds true. 

If $r=p$, using CAR for the operators $a$ gives:
\[\{ a(h_{pq}^{(\alpha,k)}), a(h_{pq}^{(\alpha',k)})\} = 0,\qquad \{ a(h_{pq}^{(\alpha,k)}), a(h_{pq}^{(\alpha',k)})^*\} = \delta_{\alpha\alpha'}1,\]
Thus
\begin{eqnarray*}
    \{(d^{(k)}_{pq}), (d^{(k)}_{pq})^*\} & = & 
    \frac{1}{4}\{a(h^{(1,k)}_{pq}) +a(h^{(1,k)}_{pq})^*
+i\,a(h^{(2,k)}_{pq})
+i\,a(h^{(2,k)}_{pq})^*    
 , \\ && \qquad 
a(h^{(1,k)}_{pq})^*
+a(h^{(1,k)}_{pq})
-i\,a(h^{(2,k)}_{pq})^*
-i\,a(h^{(2,k)}_{pq})
\} \\ &= & \frac{1}{4}(1 + 1 - (i \cdot i) - (i\cdot i)) = 1.
\end{eqnarray*}
This proves \eqref{eqn:CARd} in the remaining case.  

Rewriting \eqref{eqn:CARd} gives $d^{(k)}_{rq}(d^{(k)}_{pq})^* = \delta_{rp}1 - (d^{(k)}_{pq})^* d^{(k)}_{rq}
$, and so 
 \[d^{(k)}_{rq}(d^{(k)}_{pq})^*\Omega_{J,M} = 
\bigl(\delta_{rp}1 - (d^{(k)}_{pq})^* d^{(k)}_{rq}\bigr)\Omega_{J,M} = \delta_{rp} \Omega_{J,M}\] as claimed. 
\end{proof}



Let now 
\[
\eta_0 = \frac{1}{\sqrt{N}} \sum_{p=1}^N e_p \otimes \Omega_{J,M}
\] where $e_p$ denotes the column vector all of whose entries are zero, except that the $p$-th entry is $1$. 

\begin{lemma} \label{lemma:QOnEta} For each $k$,
    \[(Q^{(k)}_N)^* Q^{(k)}_N\eta_0 = 0,  \qquad Q^{(k)}_N (Q^{(k)}_N)^* \eta_0 = 4 \eta_0.\]
\end{lemma}

\begin{proof}
    Since $Q^{(k)}_N = 2 N^{-1/2} \sum_{p,q} e_{pq}\otimes d_{pq}^{(k)}$, we can apply Lemma~\ref{lemma:propertiesOfOmega} to deduce
    \[Q_N^{(k)} \eta_0 = \frac{2}{N}\sum_{pq} e_{p}\otimes  d_{pq}^{(k)} \Omega_{J,M} = 0  \] Thus also \( (Q^{(k)}_N)^* Q^{(k)}_N\eta_0 = 0\).

    On the other hand,
    \begin{eqnarray*}
        Q^{(k)}_N (Q^{(k)}_N)^* \eta_0 &=& 
        \frac{4}{N^{3/2}}\sum_{p,q,r} e_{r} \otimes d_{rq}^{(k)} (d_{pq}^{(k)})^* \Omega_{J,M} \\
        &=& \frac{4}{N^{3/2}}\sum_{p,q,r} e_{r} \otimes \delta_{rp}  \Omega_{J,M}  \\
        &=& \frac{4}{N^{3/2}}\sum_{q,r} e_{r} \otimes \Omega_{J,M}  
        =4\eta_0,
    \end{eqnarray*} where we used Lemma~\ref{lemma:propertiesOfOmega} to pass from the first to the second line.
\end{proof}

\subsection{Norms of sums of squares.}
Let $$R_N = \begin{bmatrix} Q_N^{(1)} & Q_N^{(2)} & \cdots & Q_N^{(M)} \\
0 & 0 & \cdots & 0 \\
\vdots & \vdots  & & \vdots \\
0 & 0 & \cdots & 0 \end{bmatrix} \in M_{M\times M}(\mathbb{C})\otimes \Cl(W,B). $$  Then $$
R_N R_N^*  = 1\otimes \sum_{j=1}^M Q_N^{(j)} (Q_N^{(j)})^*,$$ and so using Lemma~\ref{lemma:QOnEta}, $$\Vert R_N R_N^*\Vert \geq \frac{1}{\Vert \eta_0 \Vert} \sum_{j=1}^M \Vert Q_N^{(j)} (Q_N^{(j)})^* \eta_0 \Vert = \frac{4M\Vert\eta_0\Vert}{\Vert\eta_0\Vert}=4M.
$$
It follows that $$\liminf_{N\to\infty} \Vert R_N \Vert \geq 2\sqrt{M}. $$

On the other hand, let $A^{(1)},B^{(1)},\dots,A^{(M)},B^{(mM}$ be free semicircular elements, and let $Q^{(j)}=A^{(j)} + i B^{(j)}$.  be  circular elements (these are $\sqrt{2}$ times the usual normalization of a circular element).  Let $$
R = \begin{bmatrix} Q^{(1)} & Q^{(2)} & \cdots & Q^{(M)} \\
0 & 0 & \cdots & 0 \\
\vdots & \vdots  & & \vdots \\
0 & 0 & \cdots & 0 \end{bmatrix} \in M_{M\times M}(\mathbb{C})\otimes C^*(Q^{(1)},\dots,Q^{(M)}).
$$
Then $$RR^* = \sum_{j=1}^M Q^{(j)} (Q^{(j)})^*$$ is an $M$-fold sum of $\lambda=1$ free Poisson elements $Q^{(j)}(Q^{(j)})^*$ each having a law with support $[0,8]$ and density $$
\frac{1}{4\pi x}\sqrt{(8-x)x}\,
\mathbf 1_{[0,8]}(x).
$$  
Thus $RR^*$ is a $\lambda=M$  free Poisson element having its law supported on the interval 
$[2(\sqrt{N}-1)^2,2(\sqrt{M}+1)^2]$. % and having the density 
%$$
%\frac{1}{4\pi x}
%\sqrt{
%\left(2(\sqrt m+1)^2-x\right)
%\left(x-2(\sqrt m-1)^2\right)
%} \mathbf{1}_{[2(\sqrt{m}-1)^2, 2(\sqrt{m}+1)^2]}(x).
%$$ 
In particular $$\Vert R\Vert = \sqrt{2}(\sqrt{M}+1).$$

Noting that if $M\geq 6$, then $2\sqrt{M} > \sqrt{2}(\sqrt{M}+1)$, we can record the above computations as the following: 
\begin{lemma}
    \label{lemma:WrongOperatorSpaceStructure}
    With the above notation, 
    \begin{eqnarray} && \Vert R \Vert = \Big\Vert \sum_{j=1}^M Q^{(j)} (Q^{(j)})^* \Big\Vert^{1/2} = \sqrt{2}(\sqrt{M}+1),\\ && \liminf_{N\to\infty}  \Vert R_N \Vert = 
    \liminf_{N\to\infty} \Big\Vert \sum_{j=1}^M Q_N^{(j)} (Q_N^{(j)})^*\Big\Vert^{1/2} \geq 2\sqrt{M}.\end{eqnarray} xw
    In particular, if $M\geq 6$, then $2\sqrt{M}>\sqrt{2}(\sqrt{M}+1)$ and so \begin{eqnarray}
        &&\liminf_{N\to\infty} \Vert R_N\Vert \neq \Vert R \Vert,\label{eq:noOpConv}\\
        &&\liminf_{N\to\infty}\Big\Vert \sum_{j=1}^M Q_N^{(j)} (Q_N^{(j)})^*\Big\Vert \neq \Big\Vert \sum_{j=1}^M Q^{(j)} (Q^{(j)})^* \Big\Vert.\label{eq:noStrongConv}
    \end{eqnarray} 
\end{lemma}
We record an immediate corollary: 


\begin{theorem}
   \label{thrm:noStrong}
    Let $M\geq 6$.  Then: 
    \begin{itemize}
    \item[(i)] the asymptotic operator space structure of  $$\operatorname{span}_\mathbb{C}(A_N^{(1)},B_N^{(1)},\dots,A_N^{(M)},B_N^{(M)})$$  is different from that of the complex span of $2M$ free sesmicircular variables.  
    \item[(ii)] The family $(A_N^{(1)},B_N^{(1)},\dots,A_N^{(M)},B_N^{(M)})$ converges to the free semicircular family $(A^{(1)},B^{(1)},\dots,A^{(M)},B^{(M)})$ in law but not strongly. In fact, strong convergence already fails for the quadratic polynomial $$P(a_1,b_1,\dots,a_m,b_m)=\sum_{j=1}^M  (a_j + i b_j)(a_j + b_j)^*$$
    \end{itemize}
\end{theorem}

\begin{proof}
    The first part follows by  noting that $$R_N =  \sum_{j=1}^M e_{1j} \otimes A_N^{(j)} + \sum_{j=1}^M (i e_{1j})\otimes B_N^{(j)}$$ and by applying \eqref{eq:noOpConv}.

    The second part follows by evaluating $P$ and applying \eqref{eq:noStrongConv}.
\end{proof}


\end{document}
